% ============================================================================
% defensa-sus.tex
% Presentación Beamer: 3 diapositivas de resultados SUS
% ============================================================================

\documentclass[10pt,aspectratio=169]{beamer}

\usepackage[spanish,es-tabla]{babel}
\usepackage{fontspec}
\usepackage{icomma}
\usepackage{booktabs}
\usepackage{tikz}
\usepackage{pgfplots}
\pgfplotsset{compat=1.18}
\usetikzlibrary{babel}

% Tema y configuración visual de la presentación
\usetheme{Madrid}
\usecolortheme{whale}

% Tema y configuración visual de la presentación
\usetheme{Madrid}
\usecolortheme{whale}

% Carga centralizada de colores para las gráficas
% ============================================================================
% graficas/colores.tex
% Definición centralizada de colores para las gráficas SUS
% Permite modificar la apariencia de todas las gráficas desde un único lugar
% ============================================================================

\usepackage{xcolor}

% Paleta de aceptabilidad y rendimiento SUS
\definecolor{susRojo}{HTML}{D9534F}       % No aceptable (0 - 50)
\definecolor{susAmarillo}{HTML}{F0AD4E}   % Marginal (50 - 70)
\definecolor{susVerde}{HTML}{5CB85C}      % Aceptable (70 - 100)

% Colores de interfaz y visualización general
\definecolor{susPrincipal}{HTML}{1F4E79}  % Azul corporativo / barras principales
\definecolor{susSecundario}{HTML}{2E75B6} % Azul medio para detalles
\definecolor{susGrisFondo}{HTML}{E9ECEF}  % Gris claro para fondo de escalas / donut
\definecolor{susReferencia}{HTML}{C00000} % Rojo distintivo para línea de media SUS (68)
\definecolor{susTexto}{HTML}{212529}      % Texto oscuro de alto contraste


% Personalización de la paleta institucional con los colores de la plantilla
\setbeamercolor{structure}{fg=susPrincipal}
\setbeamercolor{palette primary}{bg=susPrincipal,fg=white}
\setbeamercolor{palette secondary}{bg=susSecundario,fg=white}
\setbeamercolor{palette tertiary}{bg=susPrincipal!80!black,fg=white}
\setbeamercolor{title}{bg=susPrincipal,fg=white}

% Carga del motor de cálculo en Lua, procesamiento del CSV y definición de macros
\directlua{
  sus = dofile("sus-calculos.lua")
  sus.cargar_csv("datos/respuestas.csv")
  sus.definir_macros()
}

\title[Evaluación SUS]{Evaluación de Usabilidad del Sistema (SUS)}
\subtitle{Resultados psicométricos y validación de la experiencia de usuario}
\author{Trabajo de Fin de Grado / Calidad del Software}
\date{\today}

\begin{document}

% ----------------------------------------------------------------------------
% Diapositiva de portada
% ----------------------------------------------------------------------------
\begin{frame}
  \titlepage
\end{frame}

% ----------------------------------------------------------------------------
% Diapositiva 1: Resultado global y distribución de aceptabilidad
% ----------------------------------------------------------------------------
\begin{frame}{Resultado Global de Usabilidad}
  \begin{columns}[c]
    \begin{column}{0.48\textwidth}
      \centering
      \scalebox{0.78}{% ============================================================================
% graficas/donut.tex
% Gráfica Donut: Media SUS en grande, arco proporcional, nota y adjetivo
% Puede incluirse con % ============================================================================
% graficas/donut.tex
% Gráfica Donut: Media SUS en grande, arco proporcional, nota y adjetivo
% Puede incluirse con % ============================================================================
% graficas/donut.tex
% Gráfica Donut: Media SUS en grande, arco proporcional, nota y adjetivo
% Puede incluirse con \input{graficas/donut.tex} en documentos y presentaciones
% ============================================================================

\input{graficas/colores.tex}

\begin{tikzpicture}[scale=0.9, every node/.style={transform shape}]
  % Cálculo de ángulo para el arco proporcional (0 a 360 grados)
  \edef\mediaRaw{\directlua{sus.get_stat_raw("media")}}
  \pgfmathsetmacro{\anguloMedia}{\mediaRaw * 3.6}
  \pgfmathsetmacro{\anguloFin}{90 - \anguloMedia}

  % Selección dinámica del color según aceptabilidad
  \pgfmathsetmacro{\colorDonut}{
    \mediaRaw >= 70 ? "susVerde" : (\mediaRaw >= 50 ? "susAmarillo" : "susRojo")
  }

  % Anillo de fondo (100% de la escala)
  \draw[line width=12mm, susGrisFondo] (0,0) circle (2.2cm);

  % Arco de puntuación obtenida
  \ifdim\anguloMedia pt>0.5pt
    \draw[line width=12mm, \colorDonut] (90:2.2cm) arc (90:\anguloFin:2.2cm);
  \fi

  % Textos centrales
  \node[align=center] at (0, 0.25) {
    {\fontsize{26}{30}\selectfont\bfseries \directlua{sus.get_stat("media", 1)}}\\[2pt]
    {\footnotesize\color{gray} sobre 100}
  };

  % Calificación y adjetivo Bangor / Sauro-Lewis en la parte inferior
  \node[align=center, text=susTexto] at (0, -3.2) {
    {\large\bfseries Nota: \directlua{sus.get_stat("nota_media")}}\\[3pt]
    {\normalsize\itshape \directlua{sus.get_stat("adjetivo_media")}}\\[2pt]
    {\footnotesize\bfseries (\directlua{sus.get_stat("aceptabilidad_media")})}
  };
\end{tikzpicture}
 en documentos y presentaciones
% ============================================================================

% ============================================================================
% graficas/colores.tex
% Definición centralizada de colores para las gráficas SUS
% Permite modificar la apariencia de todas las gráficas desde un único lugar
% ============================================================================

\usepackage{xcolor}

% Paleta de aceptabilidad y rendimiento SUS
\definecolor{susRojo}{HTML}{D9534F}       % No aceptable (0 - 50)
\definecolor{susAmarillo}{HTML}{F0AD4E}   % Marginal (50 - 70)
\definecolor{susVerde}{HTML}{5CB85C}      % Aceptable (70 - 100)

% Colores de interfaz y visualización general
\definecolor{susPrincipal}{HTML}{1F4E79}  % Azul corporativo / barras principales
\definecolor{susSecundario}{HTML}{2E75B6} % Azul medio para detalles
\definecolor{susGrisFondo}{HTML}{E9ECEF}  % Gris claro para fondo de escalas / donut
\definecolor{susReferencia}{HTML}{C00000} % Rojo distintivo para línea de media SUS (68)
\definecolor{susTexto}{HTML}{212529}      % Texto oscuro de alto contraste


\begin{tikzpicture}[scale=0.9, every node/.style={transform shape}]
  % Cálculo de ángulo para el arco proporcional (0 a 360 grados)
  \edef\mediaRaw{\directlua{sus.get_stat_raw("media")}}
  \pgfmathsetmacro{\anguloMedia}{\mediaRaw * 3.6}
  \pgfmathsetmacro{\anguloFin}{90 - \anguloMedia}

  % Selección dinámica del color según aceptabilidad
  \pgfmathsetmacro{\colorDonut}{
    \mediaRaw >= 70 ? "susVerde" : (\mediaRaw >= 50 ? "susAmarillo" : "susRojo")
  }

  % Anillo de fondo (100% de la escala)
  \draw[line width=12mm, susGrisFondo] (0,0) circle (2.2cm);

  % Arco de puntuación obtenida
  \ifdim\anguloMedia pt>0.5pt
    \draw[line width=12mm, \colorDonut] (90:2.2cm) arc (90:\anguloFin:2.2cm);
  \fi

  % Textos centrales
  \node[align=center] at (0, 0.25) {
    {\fontsize{26}{30}\selectfont\bfseries \directlua{sus.get_stat("media", 1)}}\\[2pt]
    {\footnotesize\color{gray} sobre 100}
  };

  % Calificación y adjetivo Bangor / Sauro-Lewis en la parte inferior
  \node[align=center, text=susTexto] at (0, -3.2) {
    {\large\bfseries Nota: \directlua{sus.get_stat("nota_media")}}\\[3pt]
    {\normalsize\itshape \directlua{sus.get_stat("adjetivo_media")}}\\[2pt]
    {\footnotesize\bfseries (\directlua{sus.get_stat("aceptabilidad_media")})}
  };
\end{tikzpicture}
 en documentos y presentaciones
% ============================================================================

% ============================================================================
% graficas/colores.tex
% Definición centralizada de colores para las gráficas SUS
% Permite modificar la apariencia de todas las gráficas desde un único lugar
% ============================================================================

\usepackage{xcolor}

% Paleta de aceptabilidad y rendimiento SUS
\definecolor{susRojo}{HTML}{D9534F}       % No aceptable (0 - 50)
\definecolor{susAmarillo}{HTML}{F0AD4E}   % Marginal (50 - 70)
\definecolor{susVerde}{HTML}{5CB85C}      % Aceptable (70 - 100)

% Colores de interfaz y visualización general
\definecolor{susPrincipal}{HTML}{1F4E79}  % Azul corporativo / barras principales
\definecolor{susSecundario}{HTML}{2E75B6} % Azul medio para detalles
\definecolor{susGrisFondo}{HTML}{E9ECEF}  % Gris claro para fondo de escalas / donut
\definecolor{susReferencia}{HTML}{C00000} % Rojo distintivo para línea de media SUS (68)
\definecolor{susTexto}{HTML}{212529}      % Texto oscuro de alto contraste


\begin{tikzpicture}[scale=0.9, every node/.style={transform shape}]
  % Cálculo de ángulo para el arco proporcional (0 a 360 grados)
  \edef\mediaRaw{\directlua{sus.get_stat_raw("media")}}
  \pgfmathsetmacro{\anguloMedia}{\mediaRaw * 3.6}
  \pgfmathsetmacro{\anguloFin}{90 - \anguloMedia}

  % Selección dinámica del color según aceptabilidad
  \pgfmathsetmacro{\colorDonut}{
    \mediaRaw >= 70 ? "susVerde" : (\mediaRaw >= 50 ? "susAmarillo" : "susRojo")
  }

  % Anillo de fondo (100% de la escala)
  \draw[line width=12mm, susGrisFondo] (0,0) circle (2.2cm);

  % Arco de puntuación obtenida
  \ifdim\anguloMedia pt>0.5pt
    \draw[line width=12mm, \colorDonut] (90:2.2cm) arc (90:\anguloFin:2.2cm);
  \fi

  % Textos centrales
  \node[align=center] at (0, 0.25) {
    {\fontsize{26}{30}\selectfont\bfseries \directlua{sus.get_stat("media", 1)}}\\[2pt]
    {\footnotesize\color{gray} sobre 100}
  };

  % Calificación y adjetivo Bangor / Sauro-Lewis en la parte inferior
  \node[align=center, text=susTexto] at (0, -3.2) {
    {\large\bfseries Nota: \directlua{sus.get_stat("nota_media")}}\\[3pt]
    {\normalsize\itshape \directlua{sus.get_stat("adjetivo_media")}}\\[2pt]
    {\footnotesize\bfseries (\directlua{sus.get_stat("aceptabilidad_media")})}
  };
\end{tikzpicture}
}
    \end{column}
    \begin{column}{0.48\textwidth}
      \centering
      \scalebox{0.78}{% ============================================================================
% graficas/velocimetro.tex
% Gráfica de velocímetro tipo tacómetro con zonas de aceptabilidad y aguja
% ============================================================================

% ============================================================================
% graficas/colores.tex
% Definición centralizada de colores para las gráficas SUS
% Permite modificar la apariencia de todas las gráficas desde un único lugar
% ============================================================================

\usepackage{xcolor}

% Paleta de aceptabilidad y rendimiento SUS
\definecolor{susRojo}{HTML}{D9534F}       % No aceptable (0 - 50)
\definecolor{susAmarillo}{HTML}{F0AD4E}   % Marginal (50 - 70)
\definecolor{susVerde}{HTML}{5CB85C}      % Aceptable (70 - 100)

% Colores de interfaz y visualización general
\definecolor{susPrincipal}{HTML}{1F4E79}  % Azul corporativo / barras principales
\definecolor{susSecundario}{HTML}{2E75B6} % Azul medio para detalles
\definecolor{susGrisFondo}{HTML}{E9ECEF}  % Gris claro para fondo de escalas / donut
\definecolor{susReferencia}{HTML}{C00000} % Rojo distintivo para línea de media SUS (68)
\definecolor{susTexto}{HTML}{212529}      % Texto oscuro de alto contraste


\begin{tikzpicture}[scale=1.1, every node/.style={transform shape}]
  \edef\mediaRaw{\directlua{sus.get_stat_raw("media")}}

  % Coordenadas angulares:
  % 0 SUS = 180 grados (izquierda)
  % 50 SUS = 90 grados (arriba)
  % 70 SUS = 54 grados
  % 100 SUS = 0 grados (derecha)
  % Fórmula de mapeo: angulo = 180 - (puntuacion * 1.8)

  % Zona No Aceptable: 0 a 50 (180 a 90 deg)
  \draw[line width=11mm, susRojo!85] (180:2.1cm) arc (180:90:2.1cm);

  % Zona Marginal: 50 a 70 (90 a 54 deg)
  \draw[line width=11mm, susAmarillo!90] (90:2.1cm) arc (90:54:2.1cm);

  % Zona Aceptable: 70 a 100 (54 a 0 deg)
  \draw[line width=11mm, susVerde!85] (54:2.1cm) arc (54:0:2.1cm);

  % Marcas y etiquetas principales de la escala
  \foreach \val in {0, 20, 40, 50, 60, 70, 80, 100} {
    \pgfmathsetmacro{\ang}{180 - (\val * 1.8)}
    \draw[white, thick] (\ang:1.55cm) -- (\ang:2.65cm);
    \node[font=\scriptsize\bfseries, text=susTexto] at (\ang:2.95cm) {\val};
  }

  % Etiquetas descriptivas de las zonas
  \node[font=\scriptsize\bfseries, text=susRojo!80!black] at (135:1.25cm) {No aceptable};
  \node[font=\tiny\bfseries, text=susAmarillo!90!black, rotate=-18] at (72:1.25cm) {Marginal};
  \node[font=\scriptsize\bfseries, text=susVerde!80!black] at (27:1.25cm) {Aceptable};

  % Cálculo del ángulo de la aguja según la media obtenida
  \pgfmathsetmacro{\angAguja}{180 - (\mediaRaw * 1.8)}

  % Dibujo de la aguja indicadora
  \draw[line width=2.4pt, susPrincipal, line cap=round] (0,0) -- (\angAguja:2.35cm);
  \fill[susPrincipal] (0,0) circle (5pt);
  \fill[white] (0,0) circle (2pt);

  % Lectura numérica y nota al pie del velocímetro
  \node[align=center, below] at (0, -0.3) {
    {\fontsize{18}{22}\selectfont\bfseries \directlua{sus.get_stat("media", 1)}}\\[2pt]
    {\small\bfseries \directlua{sus.get_stat("aceptabilidad_media")}}
  };
\end{tikzpicture}
}
    \end{column}
  \end{columns}

  \vspace{0.4cm}
  \begin{block}{Síntesis cuantitativa ($n = \susN$ participantes)}
    \centering
    Media: \textbf{\susMedia{} / 100} \quad\textbar\quad
    Calificación: \textbf{\susNotaMedia} \quad\textbar\quad
    Adjetivo: \textbf{\susAdjetivoMedia} \quad\textbar\quad
    Desv. Típica: \textbf{$s = \susDesviacion$}
  \end{block}
\end{frame}

% ----------------------------------------------------------------------------
% Diapositiva 2: Rango muestral y respuestas individuales
% ----------------------------------------------------------------------------
\begin{frame}{Dispersión e Intervalo de Puntuaciones}
  \begin{center}
    \scalebox{0.82}{% ============================================================================
% graficas/rango.tex
% Gráfica de Rango SUS: barra de 0 a 100, línea de referencia 68 y rango mín-máx
% Colores se cargan en el preámbulo principal (graficas/colores.tex)
% ============================================================================

\begin{tikzpicture}[x=0.095cm, y=1cm]
  % Usar macros definidas por sus.definir_macros() (valores con punto decimal)
  \pgfmathsetmacro{\minRaw}{\susMinimoRaw}
  \pgfmathsetmacro{\maxRaw}{\susMaximoRaw}
  \pgfmathsetmacro{\mediaRaw}{\susMediaRaw}

  % Fondo de aceptabilidad en la barra horizontal (y de 0 a 0.7)
  \fill[susRojo!25] (0,0) rectangle (50,0.7);
  \fill[susAmarillo!25] (50,0) rectangle (70,0.7);
  \fill[susVerde!25] (70,0) rectangle (100,0.7);

  % Borde exterior de la barra
  \draw[thick, gray!70] (0,0) rectangle (100,0.7);

  % Marcas de división en el eje
  \foreach \x in {0, 20, 40, 60, 80, 100} {
    \draw[gray!70] (\x, 0) -- (\x, -0.15) node[below, font=\scriptsize] {\x};
  }

  % Línea de referencia estándar de la industria (SUS = 68)
  \draw[susReferencia, very thick, dashed] (68, -0.3) -- (68, 1.3);
  \node[susReferencia, above, font=\scriptsize\bfseries, align=center] at (68, 1.35) {Media global SUS\\(68)};

  % Intervalo Mínimo - Máximo de los participantes
  \fill[susPrincipal!40, rounded corners=2pt] (\minRaw, 0.15) rectangle (\maxRaw, 0.55);
  \draw[susPrincipal, thick] (\minRaw, 0.35) -- (\maxRaw, 0.35);

  % Extremo mínimo
  \draw[susPrincipal, line width=1.2pt] (\minRaw, 0.15) -- (\minRaw, 0.55);
  \node[susPrincipal, below, font=\scriptsize\bfseries] at (\minRaw, -0.45) {Mín: \susMinimo};
  \draw[->, >=stealth, susPrincipal, thin] (\minRaw, -0.42) -- (\minRaw, 0.1);

  % Extremo máximo
  \draw[susPrincipal, line width=1.2pt] (\maxRaw, 0.15) -- (\maxRaw, 0.55);
  \node[susPrincipal, below, font=\scriptsize\bfseries] at (\maxRaw, -0.45) {Máx: \susMaximo};
  \draw[->, >=stealth, susPrincipal, thin] (\maxRaw, -0.42) -- (\maxRaw, 0.1);

  % Marcador de la media muestral obtenida
  \fill[susPrincipal] (\mediaRaw, 0.35) circle (3pt);
  \node[susPrincipal, above, font=\small\bfseries] at (\mediaRaw, 0.75) {Media: \susMedia};
  \draw[->, >=stealth, susPrincipal, thick] (\mediaRaw, 0.72) -- (\mediaRaw, 0.45);
\end{tikzpicture}
}
  \end{center}

  \vspace{0.2cm}
  \begin{center}
    \scalebox{0.72}{% ============================================================================
% graficas/barras.tex
% Gráfico de barras por participante con línea de referencia global SUS (68)
% ============================================================================

% ============================================================================
% graficas/colores.tex
% Definición centralizada de colores para las gráficas SUS
% Permite modificar la apariencia de todas las gráficas desde un único lugar
% ============================================================================

\usepackage{xcolor}

% Paleta de aceptabilidad y rendimiento SUS
\definecolor{susRojo}{HTML}{D9534F}       % No aceptable (0 - 50)
\definecolor{susAmarillo}{HTML}{F0AD4E}   % Marginal (50 - 70)
\definecolor{susVerde}{HTML}{5CB85C}      % Aceptable (70 - 100)

% Colores de interfaz y visualización general
\definecolor{susPrincipal}{HTML}{1F4E79}  % Azul corporativo / barras principales
\definecolor{susSecundario}{HTML}{2E75B6} % Azul medio para detalles
\definecolor{susGrisFondo}{HTML}{E9ECEF}  % Gris claro para fondo de escalas / donut
\definecolor{susReferencia}{HTML}{C00000} % Rojo distintivo para línea de media SUS (68)
\definecolor{susTexto}{HTML}{212529}      % Texto oscuro de alto contraste


\begin{tikzpicture}
  \begin{axis}[
    ybar,
    width=0.95\linewidth,
    height=6.2cm,
    ymin=0,
    ymax=105,
    ylabel={Puntuación SUS},
    symbolic x coords={\directlua{sus.imprimir_etiquetas_barras()}},
    xtick=data,
    nodes near coords,
    nodes near coords style={font=\footnotesize\bfseries, text=susTexto},
    bar width=18pt,
    enlarge x limits=0.15,
    ymajorgrids=true,
    grid style={dashed, gray!30},
    axis line style={gray!60},
    tick label style={font=\small},
    ylabel style={font=\small\bfseries}
  ]
    % Barras con las puntuaciones de cada participante
    \addplot[
      fill=susPrincipal!85,
      draw=susPrincipal,
      line width=0.8pt
    ] coordinates {
      \directlua{sus.imprimir_coordenadas_barras()}
    };

    % Línea de referencia estándar de la industria (SUS = 68)
    \addplot[
      color=susReferencia,
      dashed,
      line width=1.3pt,
      domain=-1:10
    ] coordinates {
      (\susPrimerId, 68) (\susUltimoId, 68)
    };
    \node[susReferencia, fill=white, fill opacity=0.85, text opacity=1, inner sep=2pt, font=\scriptsize\bfseries] 
      at (rel axis cs:0.22, 0.69) {Media estándar de referencia (68)};
  \end{axis}
\end{tikzpicture}
}
  \end{center}
\end{frame}

% ----------------------------------------------------------------------------
% Diapositiva 3: Conclusiones e interpretación de la evaluación
% ----------------------------------------------------------------------------
\begin{frame}{Conclusiones de la Evaluación SUS}
  \directlua{sus.imprimir_conclusiones_defensa()}
\end{frame}

\end{document}
